\chapter{Introdução} \label{Cap:Introdução}

O gerenciamento financeiro pessoal constitui uma das competências mais críticas e desafiadoras para os cidadãos no cenário econômico contemporâneo. Com a digitalização em massa dos meios de pagamento e a consequente substituição do dinheiro em espécie por instrumentos eletrônicos, o volume de dados transacionais gerados cotidianamente atingiu patamares sem precedentes. No entanto, o aumento na facilidade de transacionar não foi acompanhado por ferramentas automáticas eficazes de organização e inteligência orçamentária, gerando um descompasso estrutural entre o consumo e a visibilidade financeira dos indivíduos.

A popularização do cartão de crédito como o principal meio de pagamento no Brasil trouxe consigo um desafio crescente: a dificuldade de compreender e organizar os próprios gastos de forma autônoma e tempestiva. Segundo o balanço anual do setor consolidado pela \citeonline{abecs2025panorama}, as compras com cartões de crédito, débito e pré-pagos movimentaram 4,1 trilhões de reais em 2024 (alta de 10,9\% sobre 2023), dos quais 2,8 trilhões apenas no cartão de crédito, evidenciando o papel central desses instrumentos na economia nacional. O Banco Central estima mais de 200 milhões de cartões de crédito ativos no país \cite{cartacapital2025cartoes}, e o Mapa da Inadimplência da \citeonline{serasa2026mapa} aponta as dívidas com bancos e cartões de crédito como a principal causa de negativação dos consumidores brasileiros.

O crescimento no volume transacional é impulsionado tanto pelas compras em estabelecimentos físicos quanto pelo comércio eletrônico (\textit{e-commerce}): no último trimestre de 2024, os gastos com cartão de crédito cresceram 8\% sobre o mesmo período de 2023, com alta de 15\% nas compras digitais \cite{poder360_2025}. Não obstante a conveniência proporcionada por essa digitalização, a falta de padronização nas informações repassadas aos usuários nas faturas mensais impede o planejamento consciente. Pesquisa do \citeonline{spcbrasil2018controle} com a Confederação Nacional de Dirigentes Lojistas indica que 45\% dos consumidores brasileiros não controlam o próprio orçamento e que, entre os que o fazem, predominam as anotações em caderno (28\%) e as planilhas (18\%) -- práticas propensas a erros e com altas taxas de abandono ao longo do tempo.

\begin{quadro}[htb]
	\centering
	\caption{Amostra real de descrições de transações de faturas do C6 Bank e do Nubank (bloco superior) e de transcrições por OCR das mesmas faturas rasterizadas (bloco inferior), ilustrando os obstáculos tratados nesta pesquisa.}
	\label{quadro:fatura_exemplo}
	\begin{tabularx}{\textwidth}{l X}
		\toprule
		\textbf{Descrição bruta no extrato} & \textbf{Obstáculo ilustrado} \\
		\midrule
		\multicolumn{2}{@{}l@{}}{\textit{Ruído nativo das faturas (presente no arquivo emitido pela instituição)}} \\
		\texttt{jhoonymorango} & Nome informal e aglutinado de pequeno estabelecimento \\
		\texttt{DL HELPHBOMAXCOM} & Abreviação/sigla críptica de serviço digital \\
		\texttt{IFOOD *IFD*MR DELIVE} & Prefixo de intermediador de pagamento (\textit{marketplace}) \\
		\texttt{PAYPAL *BOSS LIFE} & Prefixo de \textit{gateway} internacional sobre o nome da loja \\
		\texttt{DIVIPAYPAYMENTS} & Prefixo de parcelamento aglutinado ao nome do processador \\
		\midrule
		\multicolumn{2}{@{}l@{}}{\textit{Ruído de leitura óptica (experimento de rasterização e OCR, Seção~\ref{sec:val_resultados_ocr})}} \\
		\texttt{IFD\textquotedbl{}PUNTO ITAQUA COMER} & Troca de caractere pelo OCR: asterisco do prefixo lido como aspas (verdade de referência: \texttt{IFD*PUNTO ITAQUA COMER}) \\
		\texttt{CLUBE DO INGRESSO - PARCELA 1/2} & Fusão de colunas pelo OCR (verdade de referência: \texttt{CLUBE DO INGRESSO}) \\
		\bottomrule
	\end{tabularx}
	\fonte{Elaborado pelo autor (\the\year), a partir de transações reais (nomes de estabelecimento, sem dados do titular). O ruído do bloco inferior não está presente nas faturas nativas digitais do acervo: foi produzido ao submeter as páginas rasterizadas ao Tesseract OCR na bateria comparativa da Seção~\ref{sec:val_resultados_ocr}.}
\end{quadro}

O \textbf{problema central investigado nesta pesquisa} reside na \textbf{categorização automática, precisa e escalável de transações financeiras de cartão de crédito a partir de descrições textuais curtas, heterogêneas, ruidosas e altamente despadronizadas}. Conforme ilustrado no Quadro~\ref{quadro:fatura_exemplo}, as linhas de extrato apresentam obstáculos multidimensionais:
\begin{enumerate}
	\item \textbf{Abreviações e siglas crípticas:} Nomes de fornecedores frequentemente aparecem severamente truncados, concatenados ou representados por razões sociais obsoletas (ex.: ``AMZN MKTPLACE'', ``PAG*RestSaoPaulo'', ``DL*GOOGLE'');
	\item \textbf{Prefixos de processadores e \textit{gateways} de pagamento:} Transações intermediadas por terminais eletrônicos (maquininhas) e facilitadores digitais inserem códigos e prefixos alfanuméricos indevidos (ex.: ``PG*'', ``MP*'', ``ZOOP*'', ``SUMUP*'') que ocultam a identidade real do estabelecimento;
	\item \textbf{Ruídos decorrentes de leitura óptica (OCR):} Em documentos digitalizados, escaneados ou fotografados por câmeras móveis, motores de Reconhecimento Óptico de Caracteres introduzem erros de substituição, deleção e inserção de caracteres (ex.: letras trocadas por números semelhantes), gerando palavras fora do vocabulário padrão (\textit{Out-Of-Vocabulary} -- OOV);
	\item \textbf{Ambiguidade de \textit{marketplaces} e multicategorias:} Grandes plataformas de comércio eletrônico (como Mercado Livre, Amazon e Shopee) transacionam mercadorias de tipos muito diferentes entre si (eletrônicos, vestuário, alimentação, medicamentos) sob uma mesma fatura corporativa genérica, tornando inviável a inferência da categoria apenas pelo texto bruto do canal;
	\item \textbf{Privacidade e soberania dos dados financeiros:} O tratamento de dados transacionais bancários envolve informações de alta sensibilidade pessoal, exigindo conformidade com a Lei Geral de Proteção de Dados Pessoais (LGPD -- Lei nº 13.709/2018) \cite{brasil2018lgpd}, o que restringe o envio indiscriminado de extratos a serviços terceirizados e APIs proprietárias em nuvem pública.
\end{enumerate}

A motivação deste trabalho apoia-se em duas perspectivas complementares: o avanço técnico-científico na área de Ciência da Computação e o impacto socioeconômico direto na gestão orçamentária dos cidadãos.

No âmbito da \textbf{Ciência da Computação}, o problema integra de maneira transdisciplinar a Visão Computacional/OCR, o Processamento de Linguagem Natural (PLN), os Sistemas de Recuperação de Informação Vetorial, a Engenharia de Agentes Inteligentes baseados em Grafos de Estado e as práticas de Operações de Aprendizado de Máquina (MLOps). Desenvolver e avaliar criticamente uma arquitetura que opere sob restrições severas de hardware local (sem demandar \textit{clusters} de GPU de alto custo) contribui para a literatura de Inteligência Artificial (IA) aplicada, evidenciando as possibilidades e os limites de esteiras modulares, reativas e de baixo custo computacional. As abordagens descritas na literatura (Capítulo~\ref{chp:trabalhos_relacionados}) apresentam limitações que reforçam essa necessidade: parte delas -- como a de Souza (\citeyear{souza2025extracao}) -- depende de APIs proprietárias em nuvem e de leiautes visuais catalogados, o que implica custos recorrentes, riscos à soberania dos dados financeiros e fragilidade diante de mudanças estéticas promovidas pelas emissoras; e, de modo geral, poucas tratam do ciclo de vida do modelo em produção, isto é, do monitoramento de desvios de dados e de conceito ao longo do tempo.

Sob o ponto de vista \textbf{aplicado e socioeconômico}, automatizar o processo de saneamento e classificação de despesas reduz o esforço cognitivo e a sobrecarga de tempo associados à gestão financeira pessoal. Ao visar fornecer categorias precisas, justificativas inteligíveis e um mecanismo de aprendizado contínuo com intervenção humana (\textit{Human-in-the-Loop} -- HITL), o sistema tem o potencial de ajudar os usuários a compreender seus padrões de consumo e de contribuir para a prevenção do endividamento -- benefícios condicionados ao desempenho preditivo efetivamente alcançado, que é apurado no Capítulo~\ref{Cap:Ava_Exp}.

O \textbf{objetivo geral} desta monografia é conceber e implementar uma solução computacional de ponta a ponta para a extração, normalização e classificação automática de transações financeiras em faturas de cartão de crédito, apoiada em extração com preservação de leiaute e em agentes de linguagem executados em infraestrutura local, e avaliar experimentalmente seus componentes de extração e de classificação. Essa avaliação caracteriza a acurácia preditiva alcançável e a resiliência a ruídos textuais e ópticos, sob restrições de privacidade dos dados e de sustentabilidade operacional em hardware local.

Para concretizar o objetivo geral, definiram-se os seguintes \textbf{objetivos específicos}:

\begin{enumerate}
	\item \textbf{Ingestão e extração de documentos:} desenvolver um módulo de ingestão de faturas nativas digitais com extração que preserve o leiaute, previsto arquiteturalmente também para documentos ópticos (Seção~\ref{sec:sol_extracao_hibrida});
	\item \textbf{Normalização léxica:} projetar rotinas determinísticas de remoção de prefixos de \textit{gateways} de pagamento, ruídos alfanuméricos e aliases comerciais;
	\item \textbf{Linhas de base:} implementar e medir classificadores estatísticos clássicos sobre representação textual esparsa, como referência quantitativa de comparação;
	\item \textbf{Representações densas e modelo neural:} construir uma base vetorial de busca por similaridade com \textit{embeddings} multilíngues e investigar um modelo no nível de byte;
	\item \textbf{Agente em grafo de estados:} orquestrar a decisão em um grafo que integre modelo de linguagem, ferramenta de busca e confirmação humana (HITL) em casos ambíguos;
	\item \textbf{Infraestrutura de dados em camadas e esteira MLOps:} desenvolver a infraestrutura conteinerizada da solução, com armazenamento em camadas, mensageria por eventos e rastreamento de experimentos;
	\item \textbf{Validação experimental:} avaliar quantitativamente as arquiteturas sobre faturas reais, sob teste cego com descrições de estabelecimento disjuntas entre treino e teste, confrontando as hipóteses de pesquisa com as evidências obtidas.
\end{enumerate}

O \textbf{escopo} compreende a extração de faturas em formatos digitais nativos (CSV e PDF com camada de texto; a ingestão de OFX e JSON está prevista na arquitetura, sem analisador implementado), o saneamento léxico das descrições e a classificação automática em uma taxonomia canônica de categorias orçamentárias. A extração óptica de documentos escaneados ou fotografados está implementada como componente autônomo (Tesseract OCR, Seção~\ref{sec:sol_extracao_hibrida}), não acoplado ao serviço de extração e não validado sobre material dessa natureza: a avaliação do Capítulo~\ref{Cap:Ava_Exp} opera sobre PDFs nativos digitais e arquivos CSV, lacuna registrada no Capítulo~\ref{Cap:Conclusoes}. Ficam fora do escopo a conciliação orçamentária, a detecção de fraudes, a consultoria tributária e a camada de consumo analítico (\textit{Business Intelligence}), apontadas como trabalhos futuros. Quanto aos \textbf{requisitos não funcionais}, a solução deve executar integralmente em hardware local, sem GPU em nuvem; oferecer latência compatível com o uso interativo; garantir a auditabilidade das decisões, com registro do método e da confiança de cada inferência; e explicar as predições do agente pela cadeia de raciocínio (\textit{Chain-of-Thought}) exposta em linguagem natural. Esses requisitos ancoram as escolhas do Capítulo~\ref{Cap:Metodologia}, que registra também o grau em que cada um foi atendido na versão avaliada, e são mensurados no Capítulo~\ref{Cap:Ava_Exp}; a hipótese central é formalizada na Seção~\ref{sec:sol_questoes_pesquisa} e verificada na Seção~\ref{sec:val_respostas_questoes}.

O presente documento está estruturado em seis capítulos, seis apêndices e um anexo. O Capítulo~\ref{Cap:Fund_Teo} discute os conceitos teóricos instrumentais à solução: processamento de linguagem natural, tokenização em subpalavras e bytes, representações densas e bancos de dados vetoriais, visão computacional e OCR, arquiteturas Transformer, agentes em grafos de estado e práticas de MLOps e de dados em camadas. O Capítulo~\ref{chp:trabalhos_relacionados} analisa o estado da arte por meio de um levantamento estruturado da literatura sobre saneamento pós-OCR, mineração de textos informais e extração em documentos fiscais, sintetizado em um quadro comparativo que explicita as lacunas abordadas. O Capítulo~\ref{Cap:Metodologia} detalha a solução proposta: questões de pesquisa e hipótese central, arcabouço CRISP-DM adaptado, fonte de dados e amostragem, arquitetura de dados em camadas, extração de documentos, grafo de estados do agente, salvaguardas de privacidade e estratégia de garantia de qualidade. O Capítulo~\ref{Cap:Ava_Exp} expõe o ambiente computacional, os conjuntos de dados, o protocolo de avaliação, os resultados das linhas de base e do agente híbrido, as análises de ablação e a verificação formal das hipóteses. O Capítulo~\ref{Cap:Conclusoes} sintetiza as contribuições, avalia o cumprimento dos objetivos, discute as limitações e aponta direções de continuidade. Os Apêndices~\ref{apendice:protocolo_busca} a~\ref{apendice:rotulagem} documentam, respectivamente, o protocolo do levantamento bibliográfico, o esquema do banco de dados relacional e vetorial, o repositório de código-fonte e as instruções de reprodutibilidade, o \textit{system prompt} e o contrato Pydantic do ramo LLM, as regras de normalização léxica com as listas de \textit{marketplaces} e aliases, e o dicionário de rotulagem de referência com o mapeamento das categorias do emissor; o Anexo~\ref{anexo:declaracao_ia} reproduz a declaração institucional de uso de inteligência artificial exigida pelo programa.

